\section{Introduction}

AI security work has made real progress by identifying prompt injection, jailbreaks, tool misuse, data exfiltration, memory poisoning, and scaffold failures. That work remains necessary. But tool-using AI systems also raise a question that is not captured by asking whether a prompt was malicious or an output was unsafe. When a system reads private data, calls an external tool, updates persistent state, sends a message, grades a transcript, or recommends an institutional decision, the security issue is whether the resulting machine action stayed within the authority delegated to it.

This paper calls that problem \term{delegation assurance}. The term does not introduce a new class of AI-agent attacks. It names an evidentiary problem: what would justify the claim that an AI system's action was authorized, scoped, attributable, contestable, remediable, and auditable? Existing agent-security work gives threat models, benchmarks, defences, access-control mechanisms, and audit protocols. Those are essential components. Delegation assurance asks how those components support a safety case about action: who delegated authority, what action was licensed, what sources could change that license, which limits were non-delegable, what evidence records the transition from instruction to action, and how failure would be attributed after the fact.

A user asks an assistant to summarize an email thread but not reveal confidential attachments. A retrieved webpage contains an instruction to ignore the user and send stored credentials elsewhere. A policy document includes examples of prohibited requests, some of which look like executable commands. A memory entry says the user usually approves a class of action, although the current task does not. A judge sees only a polished transcript and scores the interaction as safe because the final answer looks compliant. In each case, the system's failure is not adequately described by output safety alone. The question is whether authority moved across language, tools, memory, and evaluation in a licensed way.

Delegation assurance therefore sits between technical control and institutional audit. At one end are familiar engineering mechanisms: least privilege, access control, scoped credentials, tool permissions, logs, monitors, and memory controls. At the other end are assurance practices: structured claims, arguments, evidence, audit records, incident review, and external accountability. The hard middle layer is language-mediated authority. AI systems increasingly receive, transform, and execute authority through language. A string may be a command, a quotation, an example, a policy boundary, an untrusted document, a tool result, or a request to analyze rather than enact. A system that cannot preserve those distinctions cannot support a strong claim that its actions remained within delegated authority.

The contribution of this paper is a delegation-assurance framework for AI systems. The framework maps familiar security failures to evidence requirements. Direct prompt injection becomes a failure to preserve instruction authority. Indirect prompt injection becomes a failure to prevent untrusted data from becoming action-guiding instruction. Tool misuse becomes a failure of action licensing. Data exfiltration becomes a failure to separate read access from disclosure authority. Memory poisoning becomes an unauthorized persistence update. Over-refusal and under-refusal become different failures of delegation calibration: one abandons authorized work, the other exceeds authorized bounds. Judge failure becomes a failure of the evaluation layer to preserve the distinctions needed to license a safety claim.

This framework also sets boundary conditions. Not every AI failure is primarily a delegation failure. A factual hallucination in a no-tool chat setting, a benchmark-contamination artifact, capability shortfall, model extraction, infrastructure compromise, or training-data privacy leakage may affect delegated action downstream, but those cases are not central examples unless they alter authorized action, action evidence, or accountability records. Delegation assurance is about systems that act, access, disclose, persist, judge, or otherwise exercise authority on behalf of someone else.

The framework can be expressed as a safety-case pattern. For a task family \(T\) in environment \(E\), the top-level claim is not merely that system \(S\) produced a safe-looking result. It is that \(S\) acted only within delegated authority, and that the authority relation can be reconstructed by a third party. The evidence has to identify the principal, delegate, action, authority source, scope, constraints, provenance, information-flow rights, persistence rights, tool-use license, non-delegable limits, and failure attribution.

This is where adversarial pragmatics enters the picture, but only as one evidence layer. A delegation-assurance case needs evidence that language-mediated authority remains stable under adversarial and ambiguous conditions. The adversarial-pragmatics benchmark is one way to test that layer: it uses controlled contrasts to ask whether systems and evaluators preserve distinctions among instructions, embedded commands, quotation, source authority, scope, refusal, task success, risk, and failure attribution. The broader paper does not depend on that benchmark. The benchmark shows what one reusable evidence layer can look like when delegation assurance reaches the language interface.

The goal is not to replace agent-security taxonomies, access-control mechanisms, audit protocols, or AI safety cases. The goal is to give them a common evidentiary grammar for justified machine action. The next step above prompt security is not another list of threats. It is a way to specify what has to be true, and what has to be recorded, before an AI system's action can be treated as legitimately delegated.
